\section{From Equations to a Cover}
\label{sec:equations}

The contract of \cref{sec:traces} asks for a cover. This section constructs one
from a solved equation system. The ingredients have different origins. The
analysis interface mirrors Goblint's \texttt{Spec}~\cite{vojdani16,seidl26}.
Side-effecting equations come from Seidl et al.\ and Apinis
et al.~\cite{seidl03,apinis12}. The solver and its certificate are those of
Tilscher et al.~\cite{tilscher26}. The routing conditions, which connect the
equations to the context policy, and their discharge for each policy are ours.

\subsection{What an analysis supplies}

An analysis describes the stores at a node by an abstract state $d$. Its
meaning $\conc(d, \rho)$ also depends on $\rho$, the solved values of the
analysis's global unknowns. We suppress $\rho$ when it plays no role. The
analysis supplies the record \isaname{dg_spec}, the counterpart of Goblint's
\texttt{Spec}: one transfer per kind of edge action, an enter operation, and
two combine operations. A call needs two states. After the return, the
caller's locals come from before the call, while the globals and the result
come from the callee. Enter therefore returns a list of entry pairs $(q, e)$:
a resume state $q$ that describes the caller's store, and an entry state $e$
that describes the store in which the callee starts. Combine builds the state
after the call from $q$ and the callee's exit state.

The local soundness laws form the locale \isaname{analysis_contract}. The
concretization is monotone, every transfer loses no successor store, and
combine covers $\mathit{combine}_e(s, t)$ whenever $q$ covers $s$ and the exit
state covers $t$. In addition, enter must return a pair that covers the
caller's store and the entered store (\isaname{entry_pairs_cover}). A numeric
domain proves only its abstract operations sound. The framework derives the
transfers once for all domains, following Nipkow and Klein's generic
interpreter~\cite[Sect.~13.5]{nipkow14}. The shipped domains are Sign,
Interval, Parity, Congruence, and their reduced product.

\subsection{Context-indexed equations and routing}

For every node $v$ and context $c$ there is a local unknown $(v, c)$. Its
right-hand side joins four sources (\isaname{routed_node_rhs_parts}):
\begin{itemize}
  \item the initial state at the program entry;
  \item the transfer along each intraprocedural edge into $v$, in the same
    context $c$;
  \item the combined result of each call returning to $v$;
  \item at a procedure entry, the entry states published to it.
\end{itemize}
At a call, a \emph{routing function} chooses the callee context from the call
site, the caller context, and the entry state. Each call publishes its entry
state as a side effect to a global seed unknown for the callee and the chosen
context, which the callee's entry equation reads. A recursive call needs no
special treatment: its equation reads an exit state that depends on the
caller's own unknowns, and the solver handles the cycle as it handles a loop.

The equations choose a callee context from an abstract state, while the
concrete semantics chooses one from the concrete call through the policy
$\adm$. Two conditions make the choices agree. Both are assumptions of the
locale \isaname{routed_context}.
\begin{description}
  \item[Adequacy] (\isaname{routing_adequate}). If $\adm$ admits $c'$ for a call
    from a store that the caller's solved state describes, then enter produces
    a pair $(q, e)$ that is routed to $c'$, whose resume state covers the
    caller's store and whose entry state covers the entered store. The
    callee's entry unknown in $c'$ is solved.
  \item[Totality] (\isaname{routing_total}). $\adm$ admits at least one context
    for every call from a store that the caller's solved state describes.
\end{description}
For the unit policy and for call strings, $\adm$ admits exactly the context the
routing function computes from the call site and the caller context. Totality
is then immediate, and adequacy reduces to entry coverage and the routed entry
unknown being solved.

For entry-state contexts, the policy \isaname{routed_entry_context_rel} is
defined from the solved valuation $\sol$:
\begin{multline*}
\adm_{u,c}(e, s, s') \;=\; \{\, \mathit{route}(u, c, e') \mid {}\\
(q, e') \in \mathit{enter}(\sol(u, c)),\; s \in \conc(q),\; s' \in \conc(e') \,\}.
\end{multline*}
That is, a context is admitted when an entry pair of the solved caller state
covers both concrete stores and is routed there. The definition is not
circular. $\sol$ is fixed before the policy is applied, and the final
guarantee concerns $\NodeColl$, which no policy changes. The policy only
decides how $\NodeColl$ is split among contexts inside the proof. The
construction resembles the a posteriori soundness of Might and Manolios,
whose abstraction maps can only be built after the analysis has
run~\cite{might09}. In our case the abstraction is fixed in advance, and the
solution fixes the relation between traces and contexts. A consequence
matters for reading the result: under entry-state contexts, $\ActColl(v, c)$
is relative to the analysis that produced $\sol$.

\begin{table}[t]
\centering
\small
\begin{tabular}{@{}lp{0.6\textwidth}@{}}
\toprule
Hypothesis & Reading \\
\midrule
\texttt{analysis\_contract} & transfers, enter and combine are sound \\
\texttt{pp} (\texttt{post\_bounded}) & $\sol$ bounds every right-hand side on the solved set $V$ \\
\texttt{fwd} & an intraprocedural edge from a solved unknown with a non-empty state leads to a solved unknown \\
\texttt{comb\_fwd} & a call's continuation is solved in the caller's context \\
\texttt{routing\_adequate}, \texttt{routing\_total} & the policy agrees with the routing \\
\texttt{resolve\_sound} & callee resolution lists every procedure a covered store can call \\
\texttt{calls\_unique} & each call site has one call edge (\isaname{calls_source_unique}) \\
\texttt{is\_bot\_sound} & a state classified as bottom describes no store, so skipping its call is sound \\
\texttt{finC}, \texttt{finE} & the graph has finitely many edges \\
seed keys & seed unknowns differ from all other global unknowns \\
\bottomrule
\end{tabular}
\caption{Hypotheses of \cref{thm:routed}: the assumptions of the locales
\isaname{routed_context} and \isaname{dg_context_activation}. The theorem adds
that the program entry is solved in $c_0$ and that the initial state describes
the initial stores.}
\label{tab:hypotheses}
\end{table}

\begin{theorem}[Routed collecting soundness, \isaname{activation_collect_dg_sound}]
\label{thm:routed}
Under the hypotheses of \cref{tab:hypotheses}, if the program entry is solved
in context $c_0$, the initial state describes the initial stores under the
initial values of the global unknowns, and those values lie below the solved
ones $\rho$, then
$\ActColl(v, c) \subseteq \cover(v, c)$ for all $v$ and $c$, where
$\cover(v, c) = \conc(\sol(v, c), \rho)$ on solved unknowns and $\emptyset$
elsewhere.
\end{theorem}

The proof checks the five obligations of \cref{thm:contract} for this cover.
The initial state gives \oblig{Init}. The edge transfers and \texttt{fwd} give
\oblig{Intra}. Adequacy with entry coverage gives \oblig{Call}, adequacy with
combine soundness gives \oblig{Return}, and totality gives \oblig{Total}.
\oblig{Return} is quantified over every admitted call edge from the call site
(\cref{sec:traces}), and \texttt{calls\_unique} reduces it to the edge being
returned through.

\subsection{The solver enters through one certificate}

\Cref{thm:routed} needs a post-solution on a set of unknowns that contains the
program entry and is closed under execution steps. The verified top-down
solver~\cite{tilscher26} provides a different certificate,
\isaname{part_post_solution}, for the query $x$ it started on, its solved set
$V$, and its result $\sol$:
\snippet{part_post_solution}
\noindent Here $\sigma(\mathit{Inl}\;u)$ is the value of a local unknown $u$;
global unknowns are the $\mathit{Inr}$ side of the sum. That is, the query is
solved, $V$ is closed under the reads of right-hand sides, and $\sol$ bounds
both the right-hand side and every side effect of each solved unknown. The two certificates point in different
directions. The solver works backward from its query, the result of
\texttt{main}, while \cref{thm:routed} follows executions forward from the
entry. The bridge restricts $V$ to its \emph{live} part: the unknowns
$(u, c)$ such that $u$ belongs to a procedure $q$ whose result
$(\ctor{FunctionResult}\;q, c)$ is solved (\isaname{live_unknowns}). The
compiler ends every body with a return edge, so every node of a procedure
reaches the procedure's result through intraprocedural edges and
call-to-continuation steps. Closure under reads then puts all nodes of a
procedure into $V$ once its result is solved in that context.
\isaname{live_unknowns_cover} proves, from a well-formed program and a
returned solve, that the live set contains the program entry and is closed
forward. This is structural reachability in the graph. It does not claim
that an execution reaches the result. Three sets are involved: the solver's
solved set $V$, its live subset, on which \cref{thm:routed} is applied, and the
states the report publishes, which are taken from $V$. Publishing more
states than the proof uses is sound, because the report only has to contain
the cover.

The rest of the proof never inspects the solver. The analysis locale takes the
solver as a parameter and assumes four facts about it
(\isaname{certified_solver}): a run in its domain yields
\isaname{part_post_solution} and a finite solved set, and a run of its
executable form that returns lies in the domain and agrees with the logical
solver. These facts are proved for the vendored solver under each of its five
update rules for global unknowns (\isaname{td_certified_solver}). An update
rule decides how contributions to a global unknown accumulate: by join, by
join per contributing right-hand side, or by \emph{warrowing}, which combines
widening and narrowing~\cite{tilscher26}. Replacing the solver would not touch
the semantic layer.
